\documentclass[conference]{IEEEtran}
\IEEEoverridecommandlockouts
% The preceding line is only needed to identify funding in the first footnote. If that is unneeded, please comment it out.
%Template version as of 6/27/2024

\usepackage{cite}
\usepackage{amsmath,amssymb,amsfonts}
\usepackage{booktabs}
\usepackage{algorithmic}
\usepackage{graphicx}
\usepackage{textcomp}
\usepackage{xcolor}
\def\BibTeX{{\rm B\kern-.05em{\sc i\kern-.025em b}\kern-.08em
    T\kern-.1667em\lower.7ex\hbox{E}\kern-.125emX}}
\begin{document}

\title{Comparative Analysis of Algorithms for Unique Username Generation and Verification at Scale}

\author{\IEEEauthorblockN{1\textsuperscript{st} Rajesh Bhandari}
\IEEEauthorblockA{\textit{Department of Computer Engineering} \\
\textit{National College of Engineering}\\
Lalitpur, Nepal \\
nce080bct027@nce.edu.np}
\and
\IEEEauthorblockN{2\textsuperscript{nd} Utsav Bhattarai}
\IEEEauthorblockA{\textit{Department of Computer Engineering} \\
\textit{National College of Engineering}\\
Lalitpur, Nepal \\
nce080bct048@nce.edu.np}
}

\maketitle


\section{Introduction}
Today, we live in the 21st century, which is the age of technology development and digitization around the world. There are more than 8 billion people living all over the world and each year millions of new users start to use the Internet, and that's why today digital identity is extremely important.

With the growing popularity of online platforms, social networks, gaming services, and web applications, there has been a growing need for unique usernames. Usernames are used as public identifiers which make it possible to differentiate between users of the system. Moreover, usernames are essential for such processes as authentication, customization, social interaction and account management. At the same time, with the growing number of users of platforms, there arises a problem with generation of unique usernames due to the problem of name collision, username squatting and the lack of availability of common names.

To address these issues, various algorithms and data structures have been developed for generation and verification of usernames. Even though database indexing guarantees efficiency and accuracy, using the database alone can make the system inefficient due to an increased workload. That's why additional techniques such as Bloom Filters and distributed hashing mechanisms are used to improve scalability and performance in large-scale systems. Bloom filters are probabilistic algorithms which allow for quick membership tests and reduce database lookups by a huge margin \cite{b1,b2}. Similarly, consistent hashing allows for efficient mapping of user records to multiple servers and greatly reduces data shuffling when new servers join or leave the network \cite{b3}. These are some of the mechanisms that make it possible to deal with username verification more efficiently.

In case of username generation, there are a number of methods used to generate unique and user-friendly usernames. These include Random Suffix Generation where random numbers or alphanumeric characters are appended to the usernames in order to minimize collisions. There is also the Snowflake-based approach which helps in generating unique identifiers in a distributed computing environment and can be applied in the generation of usernames to guarantee uniqueness. Lastly, Trie-based approaches help in prefix search and generation of usernames based on user-specified prefixes such that the user is presented with usernames similar to what he/she wants.

The research focuses on conducting a comparative study of algorithms for generating unique usernames and verifying their uniqueness. Various methods will be examined in terms of their effectiveness, scalability, ability to ensure uniqueness, efficiency, and applicability in large-scale systems. The results of this research can be useful for choosing an appropriate solution to username management for contemporary web applications. 

\section{Literature Review}
Previous research related to the problem of username generation and verification can be categorized into the following three categories: probabilistic membership verification, distributed uniqueness and load balancing, and prefix-based/generative suggestions for usernames. These three categories focus on different aspects of the username management process and constitute the design space this paper compares.

\subsection{Probabilistic Membership Verification}
The first topic revolves around how a system verifies, at minimal costs and frequently, whether a given candidate string is available or not. Luo et al. offered a comprehensive literature review on Bloom filters and the multiple variations created in order to optimize their performance \cite{b1}. They started by explaining the mechanics of the Bloom Filter, in which an element is checked for membership with a small bit array, and several hash functions are used to locate that element. The authors have also talked about the Filter's main drawbacks: an adjustable false positive probability, and an increase in space needed as the set grows in size and the inability to handle deletion. In order to overcome these limitations, different
adaptations of the structure are described in the paper, including the
Counting Bloom Filters, which uses small counters instead of bits to allow deletion of elements that have become outdated and Scalable Bloom Filters, that expands the filter when new members are added to it and do not require a predetermined size of the set. The authors conclude that Bloom filters provide a good balance between memory consumption and query speed and therefore are ideal for large systems in which a quick membership verification is conducted prior to accessing the official database.

Ramabaja and Avdullahu took this line of research even further with their work on the Distributed Bloom Filter \cite{b2}. It is mostly intended to address membership scalability issues as well as decrease communication overhead of a distributed system when nodes have their own membership and there is no central node to check against. Rather than distributing the same filter to all peers, they developed a method in which nodes that communicate are assigned a unique filter for each communication pair. This enables the absence of an element to be detected by one of the peers with high probability even if the filters in themselves have a very high rate of false positives. From their experiments, it became clear that when a network of nodes uses this method, it takes only a few communication rounds to complete convergence with the set, whereas the network using identically mapped standard Bloom filters either does not converge at all or converges slowly because each node repeatedly encounters the same false positives. The above two papers show the effectiveness of Bloom Filters in the verification of usernames. This is because through them most of the usernames can be verified without accessing the main database. In addition, the distributed version of Bloom Filters can help to verify usernames across various servers and caches.

\subsection{Distributed Uniqueness and Load Distribution}
The second topic revolves around maintaining consistency and
balance within the system while spreading it across multiple machines. Karger et al. introduced Consistent Hashing as a means of distributing the data across a changing set of servers without forcing clients to have up-to-date information about active servers \cite{b3}. In comparison to traditional hashing algorithms, where addition or deletion of servers causes most of the keys to be rehashed, consistent hashing limits the ratio of keys that need to be moved on addition or deletion of a server, and provides formal proofs of that fact as well as some guarantees of load balance and fault tolerance. This feature makes consistent hashing applicable to routing username verification queries to correct shard, since it prevents the complete rehashing caused by changes in the cluster size.

Chinthareddy presented stateless Snowflake method for generating IDs which is suitable to cloud-native ecosystems \cite{b4}. The problem was without a centralized control or manually assigned worker numbers, there was no way to ensure uniqueness globally. The proposed stateless
Snowflake generates a node ID from its private IP and not from an external coordination service allowing to produce unique, roughly timestamped, IDs at high rates despite of automatic creation and destruction of instances under autoscaling. As seen from the reported results, this method allows sustaining the same throughput rate as a
traditional coordinated Snowflake installation but gets rid of the
dependence on an external leasing service. Consistent hashing solves
the task of directing a request to the right machine within a varying
cluster, whereas stateless Snowflake solves the task of creating a
unique ID on a particular machine without requesting any permission from other machines. So, they both provide means for a distributed
system to generate and direct unique IDs without coordination.

\subsection{Prefix-Based and Generative Suggestion}
The third topic revolves around how a system ensures that a name chosen by a user is unique and attractive rather than verifying one given name.
Fredkin introduced the Trie as a type of tree-like structure which allows for a really fast storage and retrieval of strings \cite{b5}. This makes Tries suitable for performing prefix searches, and since their invention, they have been used in various applications, including auto-completion, dictionaries, and recommendations. Within username management, Trie-based systems allow a system to immediately locate other possible usernames having the same required prefix. For example, if a user wants a certain name that is already occupied, a hash-based verification structure can only verify its existence but couldn't suggest similar alternatives without additional indexing mechanisms.

Bijary et al. presented Nominalist, an agentic system which can generate username suggestions with a combination of rule-based changes and inference from the large language model \cite{b6}. The proposed system consists of two agents: a creator agent and a reviewer agent. The creator agent uses a combination of pre-defined transformations, including suffix addition, dot notation, and underscore insertion, along with the output from the generative model. The reviewer agent then checks whether a username candidate is unique based on the database of already existing usernames and evaluates it on memorability, uniqueness, and typability scores. Finally, it provides the user with a ranked list of username suggestions. Unlike the Trie-based prefix suggestion, this approach aims at generating human-readable names that are culturally appropriate rather than names that are simply structurally related to the given prefix. Here, uniqueness verification is not the aim itself but the part of the process of generating and ranking potential names. This, combined with the Fredkin’s work, demonstrates that there are two distinct approaches to username suggestion that lie on the same spectrum. While a Trie provides a quick and deterministic way to find structurally similar alternatives, an agentic, LLM-based system sacrifices some latency and computation resources for more varied and context-aware username suggestions.

\subsection{Synthesis}
These three concepts address different aspects of the same problem,
with each offering a different trade-off between the dimensions that are key for username management in real-life scenarios including latency, memory footprint, collision rate, human readability, and suitability for distributed deployment. The Bloom filter and its
distributed version guarantee an extremely low latency
and a small memory overhead \cite{b1,b2}. Only a bit array and a few hash functions need to be stored in memory. However, these approaches are subject to false positives, resulting in a non-zero collision rate. Consistent Hashing does not perform a membership test but decides which node to ask when minimal disruption in case of changes in the cluster takes place \cite{b3}. The Stateless Snowflake technique generates usernames with extremely low probability of collisions and no coordination costs\cite{b4}. This makes it ideal for distributive deployment, although it outputs either a numerical or alphanumeric token instead of a human-readable name. Tries provide fast and space-efficient exact matches and prefix queries proportional to the size of the query rather than the size of the data \cite{b5}. But their memory consumption for storing a large trie can be larger than the corresponding Bloom filter for the same set. Lastly, the agentic approach based on an LLM is the only one among all five approaches that considers the human readability of usernames \cite{b6}. It can generate usernames that are more natural and user-friendly. However, it introduces additional latency and computational cost during username generation. Further, it relies on an external model or API, creating a dependency that is not present in any other approaches.

Though earlier work has been done on these methods, no existing work compares these methods directly against one another along this shared set of dimensions. The Bloom Filter and Consistent Hashing literature analyzes its contributions within the general setting of caching and set membership problems. The Stateless Snowflake system analyzes its proposal within the context of container orchestration whereas the Trie literature analyzes its structure within the general context of information storage and retrieval. The Agentic Suggestion approach focuses its evaluation on gender detection and dataset properties rather than on the systems level cost of implementing it at scale. None of these analyses were performed using the particular load profile introduced by username registration. This paper fills this gap by comparing Bloom Filters, Consistent Hashing, Snowflake-like Suffix  Generation, Trie-based Suggestion, Random Suffix generation and Agentic Username Generation using username-specific load profiles.

\section{Methodology}

\subsection{Theoretical Framework and Justification}
This paper is based on the observation made in the literature that the methods of generating and verifying usernames have been analyzed separately. Furthermore, these techniques have often been evaluated only in terms of other goals than those that are relevant to the problem of username registration. For example, Bloom filters and the technique of consistent hashing have been designed and evaluated for general-purpose caching and set representation \cite{b1,b2,b3}. Similarly, the stateless Snowflake approach has been developed to generate unique identifiers in distributed environments such as container orchestration systems \cite{b4}. The Trie data structure has been developed for general information storage and retrieval purposes \cite{b5}, while the agent-based recommendation framework has been tested on gender recognition capabilities rather than on system-level deployment costs \cite{b6}. Therefore, there is no common conceptual framework, which can be used for the comparison between the techniques, as they have been developed in various spheres of the computer science with a variety of criteria for their effectiveness.

To ensure a fair comparison, this paper relies upon a single multi-dimensional evaluation framework. The framework is based strictly on the necessities of username registration systems discovered during the comprehensive review. The evaluated criteria include query latency, memory consumption, collison rate and how easy it is for humans to remember and read generated usernames. Finally, it evaluates whether a technique is suitable for deployemnt across multiple sites. The six techniques under consideration are random suffix generation, Trie-based suggestion, Snowflake-like suffix generation, agentic username generation, Bloom filter pre-filtering, and consistent hashing. Each technique is treated as a black-box module within a username registration system. It can perform one of two functions. The first is the \textit{generation} function, which proposes a candidate username. The second is the \textit{verification/distribution} function, which determines whether the proposed username is available or not. In distributed environments, this function also identifies the node responsible for handling the verification request. The strength of this approach is based on that, in reality, username registration systems are not dependent on a single approach. Usually, they use a username generation method coupled with a verification method. Because of this, assessing both kinds under a common set of metrics makes for a more accurate comparison. Such a comparison allows one to detect not only the fastest method at a time but also the best combinations of methods for different deployment scenarios.

\subsection{Research Design}
This paper presents simulation-based experiments to compare different username generation, verification, and distribution techniques. Unlike survey or observation, the focus is on measuring the performance of algorithms rather than analyzing human behavior. In order to make comparisons unbiased, all six techniques are implemented in the same experimental environment. The data structures Bloom filter, Trie, and consistent hash ring were implemented from scratch. A simple modulo-hashing approach was also developed for comparison.

The primary variable considered in this study is the \textit{corpus scale}, defined as the number of existing usernames in the system. Three real-world datasets of different sizes (10,000, 100,000, and 1,000,000 usernames) were used throughout the experiments. To investigate how performance changes as the dataset grows, the 1,000,000 username dataset was also splitted into smaller subsets of 50,000, 100,000, and 250,000 usernames. Six corpus scales were used for evaluation of verification techniques. The original 10,000 and 100,000 username datasets were used in the experiments. In addition, subsets of 50,000, 100,000, and 250,000 usernames were extracted from the 1,000,000-username dataset. The full 1,000,000-username dataset was also included. This allowed the effect of increasing corpus size on verification performance to be studied. Distributed-deployment methods were examined using corpus scales of 5,000, 20,000, 50,000, and 100,000 keys. This was done to analyze how the remapping overhead varied with the number of keys being routed. Username generation methods were tested with 10,000, 100,000 and 1,000,000 username datasets. 

For the verification methods, several different measures were taken in each case of a different corpus sizes. In particular, there was the measure of the construction time, which is how much time it took to create the data structure, the amount of maximum memory usage when creating the structure, and the username lookup time. Username lookup latency was measured using both the average response time and the 95th-percentile response time for the candidate request set. For the Bloom filter, false-positive and false-negative rates were computed and compared against the actual dataset.

For the distributed-deployment experiment, the important factor is the proportion of the keys that are reassigned to other nodes when the cluster size increases from eight to nine nodes. Consistent hashing and modulo-hashing approaches were compared under these conditions. In addition, average and 95th-percentile latencies for routing requests to nodes were collected. For the username generation methods, four metrics were collected for each corpus size: number of attempts to produce the username, the success rate, time needed to produce the username, and readability of the generated username.

As the term “collision” has two different meanings in this study,
both of which will be utilized, it is necessary that both meanings be explicitly defined. For the purposes of the verification experiment, the
set $S$ is the set of existing user names within the corpus, while $Q$
is the stream of query strings tested against a particular structure.
The empirical false-positive rate is defined as
\begin{equation}
\text{FPR} = \frac{|\{x \in Q : x \notin S \text{ and } x \in
\widehat{S}\}|}{|\{x \in Q : x \notin S\}|},
\end{equation}
where $\widehat{S}$ represents the set of strings reported by the
structure as present, and the false-negative ratio is symmetrically
given by
\begin{equation}
\text{FNR} = \frac{|\{x \in Q : x \in S \text{ and } x \notin
\widehat{S}\}|}{|\{x \in Q : x \in S\}|}.
\end{equation}

In generation experiments, a \textit{generated-string collision} was defined differently from false positive in verification. Collision happened in case when the username obtained using generation technique turned out to be existing in the ground truth dataset after checking.
In sequential-retry experiments using three methods: 
random suffix generation, Snowflake-style suffix generation, 
and  trie-based suggestion, the rate of first-guess collision
over $N$ base name requests is defined as
\begin{equation}
\text{Collision Rate} = \frac{1}{N}\sum_{i=1}^{N} \mathbb{I}[
\text{attempts}_i > 1],
\end{equation}
where $\text{attempts}_i$ refers to the number of candidate names the algorithm has to examine before finding one which does not collide with the request $i$. The above measure cannot be applied to the agentic-style approach. As opposed to the other approaches, this one measures
a certain fixed set of candidates as opposed to generating usernames through sequential attempts. Hence, no artificial measure has been provided, and it is marked not applicable.

The success rate of the generation technique on the same $N$ requests
is expressed as 
\begin{equation}
\text{Success Rate} = \frac{1}{N}\sum_{i=1}^{N} \mathbb{I}[
\text{result}_i \neq \text{None}],
\end{equation}
which is the ratio of requests where the technique successfully produced
a username within its attempt limit, not exceeding the limit without success.

\subsection{Data Collection Procedure}
Since this study focuses on evaluating algorithms rather than human participants, no personal or identifiable user data was collected. Instead publicly available datasets with Twitter usernames were used. Three such datasets of varying sizes, with 10,000, 100,000, and 1,000,000 Twitter usernames, respectively, were chosen for testing. They were obtained from the Twitter Username Alias Dataset made available by McKelvey et al.\ \cite{b7}. It was originally assembled in support of research into aligning real names with self-selected online aliases on Twitter from Twitter user profiles. Every data file consists of one username per line, without any additional personal information or metadata beyond the username itself.

Table~\ref{tab:dataset-stats} provides the statistics about length and composition of each dataset. The mean length of usernames rises from 11.06 characters in both the 10,000 and 100,000 username datasets to 13.86 characters in the 1,000,000-username dataset, along with the standard
deviation of length (from 2.59 to 3.49) and maximum length of usernames
(from 15 to 25 characters). The proportion of usernames which include
at least one digit is roughly doubled, from 26 percent in the two smaller
datasets to 59 percent in the largest one, while the proportion of usernames containing an underscore rises almost three-fold, from 20 percent to 56 percent. The proportion of usernames which include an uppercase letter stays relatively stable, at around 52--53 percent for all three datasets.

\begin{table}[h]
\centering
\caption{Dataset length and composition statistics}
\label{tab:dataset-stats}
\begin{tabular}{lrrr}
\toprule
 & 10K & 100K & 1M \\
\midrule
Count             & 10{,}000 & 100{,}000 & 1{,}000{,}000 \\
Mean length       & 11.06    & 11.07     & 13.86 \\
Median length     & 11        & 11        & 14 \\
Std.\ dev.\ length & 2.59     & 2.58      & 3.49 \\
Min / Max length  & 2 / 15    & 2 / 15    & 2 / 25 \\
Contains digit    & 26.4\%   & 26.2\%    & 59.3\% \\
Contains `\_'     & 19.6\%   & 19.7\%    & 55.5\% \\
Contains uppercase & 52.4\%  & 53.1\%    & 53.1\% \\
Duplicates        & 0        & 0         & 0 \\
\bottomrule
\end{tabular}
\end{table}

The two datasets of 10,000 and 100,000 usernames were used diretly in the experiment. While, the 1,000,000 username dataset was used in two ways: in its entirety and as a source of sub-datasets of 50,000, 100,000, and 250,000 usernames. The latter was done by extracting the first usernames from the large dataset. The use of both the actual datasets and the sub-datasets from the largest dataset enabled the verification performance to be tested against increasing numbers of usernames.

A stream of 20,000 candidate requests was created using the 1,000,000-usernames dataset to be used for measuring latency and false positives for the verification experiment. Twenty percent of these requests consist of usernames that are direct samples from the dataset, representing the case where a user requests a username that is already taken. The other eighty percent of these requests consist of base names extracted from other entries in the dataset by removing trailing digits and separator-separated suffixes (such as underscores and periods), representing the name a user would have typed without an extra suffix. For the generation experiment, base names were extracted in the same manner from a random sample of each of the three datasets separately, so that generation performance on each scale can be measured relative to names taken from that dataset.

\subsection{Data Processing and Analysis Tools}
Every one of the six techniques was implemented as an individual Python module having a common interface. This enabled the same benchmarking setup to evaluate each technique under the same dataset and request conditions. All experiments were carried on a single machine running through one single Python process. To implement the algorithms, Python's standard library was used. While \textit{matplotlib} was used for plotting graphs and presenting the results.

The Bloom filter that was utilized in all the experiments for verification was built from scratch, where the parameters were calculated automatically based on the expected number of items to be inserted. For an expected number of elements $n$, and for a desired false-positive probability of $p = 0.01$ (1 percent), the bit array size $m$ and the number of hash functions $k$ were calculated using the formulae
\begin{equation}
m = \left\lceil \frac{-n \ln p}{(\ln 2)^2} \right\rceil, \qquad
k = \max\left(1, \text{round}\left(\frac{m}{n}\ln 2\right)\right).
\end{equation}
For example, at $n = 1{,}000{,}000$ and $p = 0.01$, this yields
$m \approx 9{,}585{,}059$ bits ($\approx 1.14$ MB) and $k = 7$ hash
functions. Instead of directly calculating $k$ independent hash functions,
the implementation adopts the double hashing method of  Kirsch and Mitzenmacher \cite{b1} to calculate all the $k$ hash values from two base hashes, MD5 and SHA-1, where $h_i(x) = (h_1(x) + i
\cdot h_2(x)) \bmod m$ for $i = 0, \dots, k-1$, thus eliminating the need for calculating seven different hash functions each time while retaining
the property of independence essential for false-positive probability
calculations.

Add-only is the policy of addition where the filter allows for insertion and membership test, but not deletion, such that once a username is inserted, it is impossible to delete it from the filter other than to re-create the entire structure, which is the limitation of the Bloom filter, which is not considered in this study, since no experiment of username deletion or account deletion was conducted. No updates to any inserted element were done at any time during the experiments; the filters were built just once per corpus size.

After loading the datasets, the benchmarking framework ran each method and automatically measured the relevant data. Various performance metrics like construction time, memory usage, query latency, collision ratio, remapping time, success ratio, and readability were captured in the experiments automatically. The query latency was measured using Python’s high-resolution performance counter, whereas memory usage was calculated using Python’s built-in memory trace utility.

Experimental results were saved in CSV files and analyzed using Python code. The analyzed data was then represented using graphs in order to analyze how the algorithms performed for varying corpus sizes. The results generated by this process were used to analyze performance trends of the methods being evaluated.



\begin{thebibliography}{00}
\bibitem{b1} L. Luo, D. Guo, R. T. B. Ma, O. Rottenstreich, and X. Luo, ``Optimizing Bloom Filter: Challenges, Solutions, and Comparisons,'' \emph{arXiv preprint arXiv:1804.04777}, 2018.

\bibitem{b2} L. Ramabaja and A. Avdullahu, ``The Distributed Bloom Filter,'' \emph{arXiv preprint arXiv:1910.07782}, 2019.

\bibitem{b3} D. Karger, E. Lehman, T. Leighton, M. Levine, D. Lewin, and R. Panigrahy, ``Consistent Hashing and Random Trees: Distributed Caching Protocols for Relieving Hot Spots on the World Wide Web,'' in \emph{Proc. 29th Annual ACM Symposium on Theory of Computing (STOC)}, El Paso, TX, USA, 1997, pp. 654--663.

\bibitem{b4} M. R. Chinthareddy, ``Stateless Snowflake: A Cloud-Agnostic Distributed ID Generator Using Network-Derived Identity,'' \textit{arXiv preprint arXiv:2512.11643}, 2025.

\bibitem{b5} E. Fredkin, ``Trie Memory,'' \textit{Communications of the ACM}, vol. 3, no. 9, pp. 490--499, Sep. 1960.

\bibitem{b6} F. Bijary, M. Ebadpour, and A. Tajbakhsh, ``Agentic Username Suggestion and Multimodal Gender Detection in Online Platforms: Introducing the PNGT-26K Dataset,'' \textit{arXiv preprint arXiv:2509.11136}, Sep. 2025.

\bibitem{b7} K. McKelvey, P. Goutzounis, S. da Cruz, and N. Chambers,
``Aligning Entity Names with Online Aliases on Twitter,'' in
\textit{Proc. 5th Int. Workshop on Natural Language Processing for
Social Media (SocialNLP@EACL 2017)}, Valencia, Spain, Apr. 2017,
pp. 25--35.

\end{thebibliography}

\end{document}
